\section{Polyhedral blocks coupled through a fixed core}\label{app:fixed-core}

This appendix proves a related elimination theorem for a single-level
problem that is linear in many small polyhedral blocks and polynomial in
a fixed-dimensional core. It gives a broader structural comparison and
makes witness recovery explicit. Its convex combinations preserve
polyhedral feasibility and linear measurements; they are not combinations
of globally optimal responses of a nonconvex follower.

Fix $r,k\ge0$ and $d\ge1$. Let $C\subset\R^r$ be compact and given by an explicit
rational quantifier-free polynomial formula. For $b=1,\ldots,B$, let
\[
 P_b(x)=\{y\in\R^{d_b}:G_b(x)y\le h_b(x)\},\qquad 1\le d_b\le d,
\]
with polynomial rational data and explicit finite rational coordinate
bounds among its rows. The local sets may be empty or lower-dimensional.
Let $A_b(x)\in\Q[x]^{k\times d_b}$, $b(x)\in\Q[x]^k$,
$c_b(x)\in\Q[x]^{d_b}$ and $c_0(x)\in\Q[x]$. Consider
\begin{equation}\label{eq:fixed-core-problem}
 \begin{aligned}
 \min_{x,z}\quad&c_0(x)+\sum_{b=1}^B c_b(x)\trans z_b,\\
 \text{subject to}\quad&x\in C,\quad z_b\in P_b(x)\quad(b=1,\ldots,B),\\
 &\sum_{b=1}^B A_b(x)z_b=b(x).
 \end{aligned}
\end{equation}
Let $L$ be the input length, including all explicit data, and let
$\delta\ge2$ bound the actual input degree. The theorem tracks the
numerical degree throughout instead of fixing it.

\begin{theorem}[Fixed-core polyhedral optimization]\label{thm:fixed-core}
For fixed $r,d,k$, feasibility and global optimization of
\eqref{eq:fixed-core-problem} are computable in time polynomial in
$(L,\delta)$. If feasible, an optimum is attained and an optimizer with
all block coordinates and its value can be returned in a single
real-algebraic field with polynomial total encoding length. Neither the
number of blocks nor their row counts is fixed. A fixed number of
aggregate inequalities is also allowed if finite bounds for their
slacks are supplied.
\end{theorem}

\subsection{Local vertices and support directions}\label{app:local-support}

For block $b$ enumerate all $d_b$-row subsets $I$. Put
\[
 \Delta_{bI}(x)=\det (G_b(x))_I,\qquad
 p_{bI}(x)=\operatorname{adj}((G_b(x))_I)(h_b(x))_I.
\]
Where $\Delta_{bI}\ne0$, the candidate is
$p_{bI}/\Delta_{bI}$. Equivalently, use the positive-denominator form
\begin{equation}\label{eq:core-vertex-formula}
 D_{bI}=\Delta_{bI}^2,\qquad Y_{bI}=\Delta_{bI}p_{bI},\qquad
                          y_{bI}=Y_{bI}/D_{bI}.
\end{equation}
Its validity predicate is
\begin{equation}\label{eq:core-vertex-validity}
 D_{bI}>0,\qquad G_b(x)Y_{bI}(x)\le h_b(x)D_{bI}(x).
\end{equation}
All these tests are polynomial of degree $O(d\delta)$.

Every vertex of $P_b(x)$ has $d_b$ independent active rows. Otherwise
some nonzero direction is orthogonal to all active normals; a small
step in either direction keeps the active rows tight and preserves the
finitely many strict inactive rows, contradicting extremality. A
nonempty bounded polyhedron has a vertex, and every linear functional
has a maximizing vertex: a maximizing face is nonempty and compact, and
successively minimizing the coordinates within it produces a singleton
face, which is a vertex of the original polyhedron. These arguments also
apply to lower-dimensional sets and singletons. Thus
\eqref{eq:core-vertex-formula} lists enough candidates to detect
nonemptiness and maximize every linear functional.

Append the objective as a measurement. With $h=k+1$, define
\begin{equation}\label{eq:core-images}
 W_b(x)=\begin{pmatrix}A_b(x)\\c_b(x)\trans\end{pmatrix},\quad
 T_b(x)=W_b(x)P_b(x),\quad
 t(x,\eta)=\begin{pmatrix}b(x)\\\eta-c_0(x)\end{pmatrix}.
\end{equation}
For a direction $\lambda\in\R^h$, the support score of a valid
vertex is
\[
 \lambda\trans W_b(x)y_{bI}(x)=n_{bI}(x,\lambda)/D_{bI}(x),\qquad
 n_{bI}=\lambda\trans W_bY_{bI}.
\]
Collect the denominator and validity polynomials, and, for every pair
of candidates in one block, the comparison polynomial
\begin{equation}\label{eq:core-score-comparison}
                     n_{bI}D_{bJ}-n_{bJ}D_{bI}.
\end{equation}
Valid candidates have positive denominators, so these signs order their
scores correctly. For fixed $d$ there are polynomially many such
polynomials, all in the $r+h$ coordinates $(x,\lambda)$, of degree
$O(d\delta)$ and polynomial coefficient bit length.

Enumerate their realizable sign conditions, including zero signs and
recording identically zero polynomials. Each sign condition determines the
valid vertices of every block and, for every nonempty block, a
support-maximizing vertex; choose the first maximizer in a fixed ordering.
Discard a condition if some block has no valid candidate. For fixed
$(r,d,k)$, the retained list $\Sigma$ has polynomial size and can be
constructed in polynomial bit time. Disconnected sign realizations cause no
problem, since signs alone determine validity and score order.

\subsection{Exact membership from the support inequalities}\label{app:support-membership}

For $\sigma\in\Sigma$, let $I_b(\sigma)$ be its selected local
index and abbreviate the corresponding denominator and numerator by
$D_{b\sigma},n_{b\sigma}$. Put
\begin{equation}\label{eq:core-cleared-support}
 \begin{split}
 H_\sigma(x)&=\prod_bD_{b\sigma}(x),\\
 J_\sigma(x,\eta,\lambda)&=
 \lambda\trans t(x,\eta)H_\sigma(x)
             -\sum_b n_{b\sigma}(x,\lambda)\prod_{a\ne b}D_{a\sigma}(x).
 \end{split}
\end{equation}
On that sign condition $H_\sigma>0$, and $J_\sigma\le0$ states the
support inequality
\[
 \lambda\trans t(x,\eta)\le
                     \sum_b\max_{y\in P_b(x)}\lambda\trans W_b(x)y.
\]
Writing $S_\sigma(x,\lambda)$ for the sign conjunction, define
\begin{equation}\label{eq:core-support-predicate}
 \Phi(x,\eta,\lambda)=\bigvee_{\sigma\in\Sigma}
                  [S_\sigma(x,\lambda)\wedge J_\sigma(x,\eta,\lambda)\le0].
\end{equation}

\begin{lemma}[Support membership]\label{lem:core-support-membership}
For $x\in C$, there exist blocks satisfying
\eqref{eq:fixed-core-problem} with objective exactly $\eta$ if and only if
\begin{equation}\label{eq:core-membership}
                           \forall\lambda\in\R^h\ \Phi(x,\eta,\lambda).
\end{equation}
\end{lemma}
\begin{proof}
If a block is empty, then for every $\lambda$ no retained sign condition
is realized at $(x,\lambda)$, so \eqref{eq:core-membership} is false.
Otherwise the Minkowski sum
$T(x)=\sum_bT_b(x)=\{\sum_b t_b:t_b\in T_b(x)\}$ is compact and
convex. Its support function is the sum of the support functions of its
summands, since a maximizer can be chosen independently in every block.
Hence \eqref{eq:core-support-predicate} gives its exact support
inequality in every direction.

A point $t$ in a compact convex set $T$ satisfies all its support
inequalities. Conversely, if $t\notin T$, choose a nearest point
$p\in T$ and put $\lambda=t-p\ne0$. The derivative of squared
distance along the segment from $p$ to any $u\in T$ gives
$(t-p)\trans(u-p)\le0$. Hence
$\max_{u\in T}\lambda\trans u\le\lambda\trans p
<\lambda\trans t$, violating a support inequality.
Thus \eqref{eq:core-membership} is equivalent to $t(x,\eta)\in T(x)$.
By \eqref{eq:core-images}, this membership states precisely existence
of blocks satisfying all linking equations and attaining the specified
objective value.
\end{proof}

The degrees in \eqref{eq:core-cleared-support} are
$O((B+1)d\delta)$, which need not be constant. Their ambient dimension
$r+h+1$ is fixed, however, so by the product argument in
\cref{lem:candidate-size} their expanded monomial count and bit size
are polynomial in $(L,\delta)$. Define
\[
 E(x,\eta)=(x\in C)\wedge\forall\lambda\,\Phi(x,\eta,\lambda).
\]
This formula in a fixed number of variables describes exactly the
attainable core--value pairs. The feasible set of
\eqref{eq:fixed-core-problem} is closed and lies in the product of the
compact set $C$ and the finite block boxes, so the minimum is attained
whenever the problem is feasible. Apply quantifier elimination and joint
sampling to
\begin{equation}\label{eq:core-joint-optimum}
 E(x,\eta)\wedge
       \neg\exists x',\eta'\,[E(x',\eta')\wedge\eta'<\eta].
\end{equation}
This gives an optimal $(x^*,\eta^*)$ in one field
$K=\Q(\alpha)$ of polynomial degree and encoding length.
It remains to recover all block variables. Making each of them a
quantified coordinate would destroy the fixed-dimensional argument.

\subsection{Constructive recovery from a few support tuples}\label{app:core-recovery}

The support enumeration already supplies a polynomial-size set of
full block tuples that suffices for recovery. For each $\sigma\in\Sigma$,
evaluate its selected local denominators at $x^*$ and discard the tuple
if any is zero. Otherwise evaluate the vertex formulas and retain the
tuple if every selected local vertex is feasible at $x^*$. Denote the
retained tuples by $y^1,\ldots,y^M$, and set
\[
              a^j=\sum_b W_b(x^*)y_b^j\in K^h\qquad(j=1,\ldots,M).
\]
The filter does not require the sign condition $\sigma$ to be realizable
at $x^*$ with some direction: feasibility of the tuple suffices, and the filter
never adds a point outside $T(x^*)$.

\begin{lemma}[A small common convex combination recovers every block]
\label{lem:core-recovery}
The retained aggregate points satisfy
$\conv\{a^1,\ldots,a^M\}=T(x^*)$. For every target $t\in K^h\cap T(x^*)$
one can find at most $h+1$ retained tuples and nonnegative weights in $K$
that sum to one and whose weighted aggregate is $t$, in polynomial bit
time in the original input and target encoding. Using the same weights
in every block gives a feasible original tuple with aggregate $t$.
\end{lemma}
\begin{proof}
Every retained tuple is blockwise feasible, so its aggregate belongs to
$T(x^*)$ and the indicated convex hull is contained in $T(x^*)$.
For every real direction $\lambda$, the actual sign condition at
$(x^*,\lambda)$ is retained in $\Sigma$: no block is empty at an
optimal core. Its selected vertices give a feasible tuple that attains
the support of $T(x^*)$ in direction $\lambda$. This tuple survives
the filter. Hence the convex hull has the same support function as
$T(x^*)$. The separation argument in \cref{lem:core-support-membership}
proves equality.

The usual Carath\'eodory reduction gives a convex representation of
a target in a finite convex hull with affinely independent support.
To see this directly, start with a representation whose weights on its
support are all positive. If its lifted
columns $(a^j,1)$ are dependent, choose a nonzero dependence $u$.
Its components sum to zero, so it has a positive component. Replace
the weights $\theta_j$ by $\theta_j-su_j$, with
$s=\min_{u_j>0}\theta_j/u_j$. This preserves the target and the sum
of weights, keeps weights nonnegative, and removes at least one support
point. Repeat until the lifted columns are independent. Their number
is at most $h+1$.

For construction, enumerate all subsets of at most $h+1$ retained
points. On each affinely independent subset solve
\[
          \sum_j\theta_j a^j=t,\qquad \sum_j\theta_j=1.
\]
Choose independent rows to solve this fixed-size system, then check
all remaining equations and the nonnegativity of its unique weights.
If a suitable subset exists over the reals, the unique solution of its
independent linear system lies in $K$, since every coefficient and the
target belong to $K$. Thus the enumeration finds it. There are at most
$\sum_{q=1}^{h+1}\binom Mq$ subsets, polynomial for fixed $h$.
All arithmetic, zero tests, and comparisons take place in the already
sampled ordered field $K$, whose degree and coefficient encodings are
polynomial. Each system has fixed order, so determinant expressions
and rational evaluation give polynomial coefficient lengths for the
weights; summing the resulting fixed number of terms preserves the
same bound for every recovered block coordinate.

Finally, put $z_b=\sum_j\theta_j y_b^j$ for every block. Convexity
of $P_b(x^*)$ makes $z_b$ feasible, and linearity of $W_b(x^*)$ gives
$\sum_bW_b(x^*)z_b=t$. The same weights must be used in all blocks:
this is how the prescribed total measurements are preserved.
\end{proof}

Applying \cref{lem:core-recovery} to
$t=t(x^*,\eta^*)$ recovers the original optimizer in the field $K$,
which proves the main claim of \cref{thm:fixed-core}. With no blocks,
the image sum is $\{0\}$ and the empty tuple is its only support tuple;
empty sums and products have their usual meanings. Blocks of dimension
zero can be removed after testing their scalar rows. A fixed number of
aggregate inequalities become equalities with bounded scalar slack
blocks; the supplied slack bounds keep the model compact and do not
change the fixed structural dimensions.

The running time is polynomial in the numerical input degree as well as
the input length: local determinants and pair comparisons have degree
$O(d\delta)$; the global products have degree $O((B+1)d\delta)$; sign
enumeration, elimination and sampling are all in fixed dimension; and
reconstruction adds only fixed-size algebraic linear systems. Under
polynomially bounded degree, the time is polynomial in $L$.

The common support-direction construction has a direct predecessor in
Gritzmann and Sturmfels~\cite[Algorithm 2.3.6, Theorem 2.3.7, and
Corollary 2.3.10]{GritzmannSturmfels1993}: for fixed-dimensional input
polytopes, they enumerate a shared arrangement of support directions and
compatible maximizing summand vertices, with polynomial binary complexity.
Beyond that construction, we prove that the sign enumeration is uniform
over the polynomially varying core, and that exact optimization and
reconstruction recover all original blocks in one common field. The
support characterization, finite-dimensional convex-hull reduction and
real algebraic algorithms are established tools. An alternative recovery
would solve a linear program over the sampled common field; Adler and
Beling~\cite{AdlerBeling1994} analyze that route with explicit dependence
on the common extension degree. The recovery above instead reuses the
support tuples already generated and needs no large-dimensional algebraic
LP solve. It relies on the affine block measurements, including the
objective coordinate, and would not justify convexifying a nonconvex
follower response set or mixing its tied responses to create additional
follower choices.
